\documentclass[tikz,border=10pt]{standalone}
\usepackage{tikz}
\usetikzlibrary{shapes.geometric, arrows.meta, positioning, fit, backgrounds, calc}
\usepackage{xeCJK}
\setCJKmainfont{SimSun}

\begin{document}
\begin{tikzpicture}[
    node distance=0.8cm and 1.5cm,
    box/.style={rectangle, rounded corners, draw=black, thick, minimum width=3.5cm, minimum height=0.8cm, align=center, font=\small},
    mainbox/.style={box, fill=blue!30, minimum width=5cm, minimum height=1cm},
    ch1box/.style={box, fill=gray!20, minimum width=2.8cm},
    ch2box/.style={box, fill=gray!20, minimum width=2.8cm},
    ch3box/.style={box, fill=green!30},
    ch4box/.style={box, fill=orange!30},
    ch5box/.style={box, fill=purple!30, minimum width=5cm},
    ch6box/.style={box, fill=gray!40, minimum width=5cm},
    methodbox/.style={box, fill=white, minimum width=2.5cm, minimum height=0.7cm},
    titlebox/.style={box, font=\small\bfseries, minimum height=0.7cm},
    dashedgroup/.style={draw=black, dashed, rounded corners, thick, inner sep=8pt},
    arrow/.style={-Stealth, thick}
]

% Chapter 1 and 2 (context)
\node[ch1box] (ch1) at (-4, 2) {第一章：绪论};
\node[ch2box] (ch2) [below=0.5cm of ch1] {第二章：相关背景知识};

% Main problem
\node[mainbox] (problem) at (0, 0) {分布式集群多节点\\时序异常检测};

% Chapter 3 branch (left)
\node[titlebox, ch3box] (ch3title) at (-3, -2.5) {第三章：单指标多节点异常检测};

\node[methodbox] (method1) at (-4.5, -4) {欧氏距离};
\node[methodbox] (method2) at (-3, -4) {自编码器嵌入};
\node[methodbox] (method3) at (-1.5, -4) {DBSCAN聚类};

\node[methodbox] (ispot) at (-3, -5.2) {I-SPOT自适应阈值};
\node[methodbox] (ensemble) at (-3, -6.4) {保守集成决策};

% Chapter 4 branch (right)
\node[titlebox, ch4box] (ch4title) at (3, -2.5) {第四章：多指标多节点异常检测};

\node[methodbox] (numeric) at (2.2, -4) {数值视角\\（欧氏距离）};
\node[methodbox] (shape) at (3.8, -4) {形状视角\\（FR-SAX）};

\node[methodbox] (weighted) at (3, -5.5) {加权融合决策};

% Chapter 5
\node[ch5box] (ch5) at (0, -8.5) {第五章：异常检测系统\\设计与实现};

% Chapter 6
\node[ch6box] (ch6) [below=0.8cm of ch5] {第六章：总结与展望};

% Dashed grouping for Chapter 3
\begin{scope}[on background layer]
\node[dashedgroup, fit=(method1)(method2)(method3)(ispot)(ensemble), fill=green!5] (ch3group) {};
\end{scope}

% Dashed grouping for Chapter 4
\begin{scope}[on background layer]
\node[dashedgroup, fit=(numeric)(shape)(weighted), fill=orange!5] (ch4group) {};
\end{scope}

% Arrows from Ch1 and Ch2 to problem
\draw[arrow] (ch1.east) -| ($(problem.north west)+(0.5,0)$);
\draw[arrow] (ch2.east) -| ($(problem.north west)+(0.5,0)$);

% Arrow from problem to Ch3 and Ch4 titles
\draw[arrow] (problem.south) -- ++(0, -0.5) -| (ch3title.north);
\draw[arrow] (problem.south) -- ++(0, -0.5) -| (ch4title.north);

% Arrows within Ch3
\draw[arrow] (method1.south) |- (ispot.west);
\draw[arrow] (method2.south) -- (ispot.north);
\draw[arrow] (method3.south) |- (ispot.east);
\draw[arrow] (ispot.south) -- (ensemble.north);

% Arrows within Ch4
\draw[arrow] (numeric.south) |- (weighted.west);
\draw[arrow] (shape.south) |- (weighted.east);

% Arrows from Ch3 and Ch4 to Ch5
\draw[arrow] (ensemble.south) |- ($(ch5.north)+(-1,0.5)$) -- ($(ch5.north)+(-1,0)$);
\draw[arrow] (weighted.south) |- ($(ch5.north)+(1,0.5)$) -- ($(ch5.north)+(1,0)$);

% Arrow from Ch5 to Ch6
\draw[arrow] (ch5.south) -- (ch6.north);

\end{tikzpicture}
\end{document}
